% part: intuitionistic-logic
% chapter: introduction
% section: syntax

\documentclass[../../../include/open-logic-section]{subfiles}

\begin{document}

\olfileid{int}{int}{syn}

\olsection{لغة المنطق الحدسي وتركيبها}

تركيب المنطق الحدسي هو تركيب المنطق القضي نفسه، غير أن الروابط لا يختصر
بعضها ببعض كما يصنع في المنطق الكلاسيكي. ففي الكلاسيكي يحد~$!A \lif !B$
بــ$\lnot !A \lor !B$، ويحد~$!A \lor !B$ بــ
$\lnot(\lnot !A \land \lnot !B)$، فتجعل بعض الكتب هذه الروابط اختصارات.
ولا تصح هذه المساواة في الحدسي إلا في موضعين: فقد يعرّف~$\lnot !A$، وهو
الغالب، اختصارًا لــ$!A \lif \lfalse$؛ وعند ذلك لا بد أن تكون~$\lfalse$
أولية غير معرّفة!{} ويعرّف كذلك~$!A \liff !B$، كما في الكلاسيكي، بــ
$(!A \lif !B) \land (!B \lif !A)$.

وتُبنى !!^{formula}s المنطق القضي الحدسي من
\emph{!!{propositional variable}s} والثابت القضي~$\lfalse$ باستعمال
\emph{الروابط المنطقية}. ولدينا:

\begin{enumerate}
\item مجموعة~$\PVar$ !!^a{denumerable} من ال!!{propositional variable}s
  $\Obj p_0$، $\Obj p_1$، \dots
  \item الثابت القضي لل!!{falsity}~$\lfalse$.
\item الروابط المنطقية: $\land$
  (العطف)، و$\lor$ (الفصل)، و$\lif$ (الشرط)
\item علامات الترقيم: (، )، والفاصلة.
\end{enumerate}

\begin{defn}[الصيغة]
\ollabel{defn:formulas} تُعرَّف المجموعة~$\Frm[L_0]$ المؤلفة من
\emph{!!{formula}s} المنطق القضي الحدسي استقرائيًا كما يأتي:
\begin{enumerate}
\item $\lfalse$ !!{formula} ذرية.

\item كل !!{propositional variable}~$\Obj p_i$ هي !!{formula} ذرية.

\item إذا كانت $!A$ و$!B$ !!{formula}s، فإن $(!A \land
  !B)$ !!a{formula}.

\item إذا كانت $!A$ و$!B$ !!{formula}s، فإن $(!A \lor !B)$
  !!a{formula}.

\item إذا كانت $!A$ و$!B$ !!{formula}s، فإن $(!A \lif !B)$
  !!a{formula}.

\tagitem{limitClause}{ولا شيء غير ذلك !!a{formula}.}{}
\end{enumerate}
\end{defn}

ومع الروابط الأولية المتقدمة نستعمل الرمزين
\emph{المعرّفين} الآتيين: $\lnot$ (النفي) و$\liff$
(!!{biconditional}). وتُفهم الصيغ المبنية باستعمال المؤثرات المعرّفة
كما يأتي:
\begin{enumerate}
\item $\lnot !A$ اختصار لـ$!A \lif \lfalse$.

\item $!A \liff !B$ اختصار لـ$(!A \lif !B) \land (!B
  \lif !A)$.
\end{enumerate}

ومع أن~$\lnot$ اختصار من جهة الرسم الصوري، فقد نفرد له أحيانًا قواعد
وبنودًا لــ$\lnot$ كما لو كان أوليًا؛ وإنما نفعل ذلك لتيسير صوغ المسائل
التدريبية.

\end{document}
